% ⛔철회된 레거시 코퍼스 유래 Discussion 절 --- **이동 보관**(2026-09-01)
% ⛔`\input` 하지 마라. 인용도 하지 마라. 보존 목적이다.
% 사유 = retracted_results_100ep.tex 헤더와 같다.

% ── (1) §6.2 Directional deficit · §6.3 BF2 failure case ──
\subsection{Directional estimation deficit (supplementary analysis)}
\label{sec:dirdeficit}

In addition to the five-metric panel, we examined a supplementary
metric---the $R^2$ between predicted and ground-truth 2D wavenumber
spectra (computed from the FFT of the spatial field on the same
evaluation grid for all methods)---to assess directional structure.
This metric is sensitive to wavenumber-bin alignment and can yield
strongly negative values when sparse spectral peaks are misaligned by
even a single bin, so we report it as a diagnostic rather than as
a primary panel item.  Under the corrected tracker WaveSplat attains
$+0.381$ across the full 100-episode dataset versus $+0.059$ for ETKF---%
a sign reversal from the $-0.03$ reported before the correction, and now
the best value among the methods compared.
The aggregate remains regime-split rather than uniform: $+0.983$ for
$\leq$5-component episodes against $-0.522$ for $\geq$15-component ones,
where single-bin peak misalignment drives the statistic strongly negative.
Direction information is therefore captured accurately in the sparse
regime and lost in the crowded one, rather than being uniformly weaker
than amplitude or shape as the earlier aggregate suggested.
A physical explanation is that forward-looking shipboard LiDAR
illuminates a narrow angular sector ahead of the vessel: the
observation swath spans $\sim$250~m in range but only
${\sim}120^\circ$ in azimuth, and the raydrop geometry preferentially
keeps returns along the heading axis.
This geometric bias means that the beamforming power $P(k_x, k_y)$
integrates energy predominantly along the heading direction, making
the along-heading wavenumber component much better resolved than the
cross-heading component.  A wave whose energy is correctly captured in
terms of $|k|$ and $\omega$ but whose propagation direction is
ambiguous between, say, $\theta$ and $-\theta$ relative to the vessel
heading will score zero on Pearson direction correlation while
achieving good nRMSE.  This is the sensor geometry analog of the
phase ambiguity known in radar wave spectroscopy~\citep{nieto2004}.

Two complementary mitigations exist: (a) exploiting vessel yaw
variation to sample multiple heading angles during initialization,
resolving the $\pm\theta$ ambiguity; or (b) incorporating a
prior over the dominant wave direction from onboard meteorological
data or ship motion analysis.  Both are deferred to future work;
users of WaveSplat for directional-sensitive downstream applications
(e.g., heading-dependent ride control) should be aware of this
limitation.

\subsection{Failure case: BF2 1-component}
\label{sec:bf2case}

On the originally released dataset WaveSplat's worst case occurred on the
simplest episode: BF2, $H_s$=0.3~m, a single component.  It is no longer
the worst case---under the Nyquist-adequate render and the corrected
tracker it ranks 45th of the 100 episodes, and the worst cases are now the
30-component seas as one would expect.  We nevertheless report it at
length, because our own initial diagnosis of it was wrong twice over, and
each correction changes what a practitioner should take from it.

On the originally released dataset this episode gave nRMSE~$=0.315$
with adaptive $N$ selecting $N$=11 for a 1-component swell, and we
attributed the failure to an intrinsically low signal-to-noise
ratio---wave amplitudes of order the LiDAR noise floor
($\sim$0.02~m)---and proposed an $H_s$-tuned regularization schedule
as the remedy.  That attribution does not survive scrutiny.  The
episode's dominant wavelength is $\lambda \approx 8.2$~m, rendered on
a raycast mesh of $6.9$~m spacing: $1.2$ points per wavelength, below
the Nyquist rate of $2$.  The wave was therefore not weakly
represented in the reference field, it was \emph{not representable in
it at all}, and the estimator was being scored against a reference
that could not contain the signal it was asked to find.  This is a
property of the simulator's rendering, not of the sea state.

Re-running the identical episode, sea state, and estimator on a
Nyquist-adequate re-rendering (Section~\ref{sec:regen}) moved the
median over five estimator seeds from $0.315$ to $0.169$, and the
subsequent tracker phase-advance correction moved it to
$\mathbf{0.047}$---a factor of $6.7$ in total, entirely from the
evaluation and tracking pipeline rather than from any change to the
estimator's design.  The $H_s$-tuned regularization schedule we
originally proposed addresses a problem that does not exist.

The re-run does not, however, exonerate the model-order mechanism, and on
this episode the corrected data make the coupling unusually legible.
Across the five seeds the final model order ranges from $N$=1 to $N$=15 on
this single fixed episode, and nRMSE increases monotonically with it:
$N{=}1 \to 0.023$, $N{=}3 \to 0.030$, and the three $N{=}15$ seeds giving
$0.047$, $0.061$, $0.071$.  Variance ratio falls along the same ordering
($0.84 \to 0.84 \to 0.68$--$0.77$), i.e.\ the over-ordered fits split one
swell's energy across spurious components.  Spurious growth on calm seas
is therefore real, and it is a \emph{convergence-basin} phenomenon of a
non-convex fit with random phase initialization---see
Section~\ref{sec:limitations}---rather than a sea-state-dependent
signal-to-noise limit.  The practical implication changes accordingly:
what helps is multiple restarts with median reporting, not an
$H_s$-tuned regularizer.  This is also the clearest single-episode view of
why the model-order mechanism carries no aggregate accuracy benefit
(Section~\ref{sec:adaptive}): its seed-to-seed variance on a
1-component sea is larger than its mean effect over 12 episodes.

We report this at length because the error class generalizes.  A
rendering limitation of the evaluation pipeline is easy to mistake for
a physical or algorithmic limit, since both produce a large error on a
case where the physics is genuinely marginal.  Wherever this paper
attributes a failure to physics or to the estimator, the attribution
has been re-checked against the mesh adequacy of the episodes
involved.


% ── (2) §6.7 Strong baselines · §6.8 Spectral accuracy · §6.9 Sparse-LETKF ──
\subsection{Comparison with stronger baselines}
\label{sec:strongbaselines}

\emph{Provenance note.}  All quantitative tables in this subsection---the
stronger-baseline comparison (Table~\ref{tab:strong_baseline}), the
paired-run Wilcoxon test (Section~\ref{subsec:stats}), and the
spectral-accuracy table (Table~\ref{tab:spectral_accuracy})---have been
recomputed on the Nyquist-adequate regenerated dataset
(Section~\ref{sec:regen}); the WaveSplat arm is the per-episode median
over 5 estimator seeds and the grid and GPR baselines follow their
original single-filter-seed protocol.  Regeneration \emph{widens} the
headline: the localized LETKF baseline is \emph{mesh-invariant}
(its nRMSE moves by $-0.0006$, paired $t=-1.56$ over 12 episodes,
because its spectral forecast is a per-mode rotation that does not pass
through the raycast mesh), so WaveSplat's improvement to $0.115$ against
% RETRACTED 2026-08-14 (one number): "0.319" is on the citation-ban list -- it is the
% ETKF column of results/b1_full_comparison.json, in which WaveSplat was scored on
% 120-220 frames while ETKF/GPR were scored on 301, in 100/100 episodes.
% The mesh-invariance observation itself is unaffected in direction; only the anchor
% value is withdrawn. (basis: memory/reference_citation_bans.md (08-13 round 3))
LETKF's steady $0.228$ is a genuine gain; the tuned ETKF is likewise
near-mesh-invariant in nRMSE (moving to $0.332$ from a pre-regeneration
value that is \textbf{[RETRACTED 2026-08-14]} --- that anchor comes from a
comparison in which the arms were not scored on the same frames),
its variance ratio
rising toward unity on the adequate reference field while its correlation
with truth does not.  Two supplementary probes below---the
localization-radius sensitivity sweep and the Beaufort~6--7 scalability
check---retain the original single-seed evaluation and, being LETKF-based
and mesh-invariant, are unaffected.

To stress-test WaveSplat's advantage beyond the ETKF and single-frame
GPR used in the primary evaluation, we implemented two stronger
baselines: (1)~a \emph{Local Ensemble Transform Kalman Filter}
(LETKF) with Gaspari--Cohn localization at radius $r_{\text{loc}}=50$~m
\citep{hunt2007}, and (2)~a \emph{dispersion-aware GPR} that
phase-corrects observations from the previous three frames using the
estimated dominant wavenumber before spatial regression.  Both methods
address known weaknesses of their simpler counterparts: LETKF prevents
the spurious long-range covariance contamination that degrades vanilla
ETKF, while multi-frame GPR incorporates temporal coherence via
wave physics.

To validate this comparison at full scale, we ran LETKF on all 100
episodes.  Table~\ref{tab:strong_baseline} reports the WaveSplat
vs.\ LETKF head-to-head on the complete dataset, alongside the
12-episode comparison that includes all five methods.

\begin{table*}[t]
\centering
\caption{WaveSplat vs.\ LETKF on the full 100-episode dataset (top)
and all five methods on 12 representative episodes (bottom).
Accuracy metrics (nRMSE, VR, $\rho$, SS, Dir.\ $R^2$, $H_s$ err) are
re-measured on the Nyquist-adequate regenerated dataset, WaveSplat as
the median over 5 estimator seeds; per-frame times (ms) are carried from
the original measurement (timing is independent of the reference mesh).
LETKF uses Gaspari--Cohn localization ($r_{\text{loc}}=50$~m,
$N_e=100$).  Bold indicates best per column.
Bootstrap 95\% CI for SS (100~ep): WaveSplat $[{+0.42},\,{+0.65}]$,
LETKF $[{+0.13},\,{+0.16}]$.}
\label{tab:strong_baseline}
\begin{tabular}{lccccccr}
\toprule
Method & nRMSE $\downarrow$ & VR $\to$ 1 & $\rho$ $\uparrow$ &
SS $\uparrow$ & Dir.\ $R^2$ $\uparrow$ & $H_s$ err $\downarrow$ &
ms $\downarrow$ \\
\midrule
\multicolumn{8}{l}{\emph{Full 100-episode comparison}} \\
\textbf{WaveSplat} & \textbf{0.115} & \textbf{0.916} & \textbf{0.764} &
\textbf{+0.538} & \textbf{+0.381} & \textbf{0.102} & \textbf{14.5} \\
LETKF (loc50)      & 0.228 & 0.127 & 0.357 & +0.144 &
+0.094 & 0.667 & 118.3 \\
\midrule
\multicolumn{8}{l}{\emph{12-episode subset (incl.\ ETKF, GPR)}} \\
\textbf{WaveSplat} & \textbf{0.182} & \textbf{0.773} & \textbf{0.580} &
\textbf{+0.310} & \textbf{+0.458} & \textbf{0.185} & \textbf{13.2} \\
LETKF (loc50)      & 0.237 & 0.129 & 0.345 & +0.131 & 0.080 & 0.646 & 115.7 \\
ETKF (tuned)       & 0.377 & 1.431 & 0.091 & $-$1.227 & 0.029 & 0.284 & 103.9 \\
GPR-1              & 0.274 & 0.046 & 0.170 & $-$0.175 & 0.013 & 0.801 & 208.8 \\
GPR-disp (3f)      & 0.268 & 0.073 & 0.224 & $-$0.119 & 0.028 & 0.746 & 473.6 \\
\bottomrule
\end{tabular}
\end{table*}

Three findings merit attention.

\paragraph{Regime-dependent advantage.}
On the full 100-episode dataset, WaveSplat wins 58/100 episodes on
nRMSE and 58/100 on Skill Score---a majority on both.
This aggregate masks a strong regime dependence: for $\leq$5 spectral
components (60~episodes), WaveSplat achieves mean SS~$=+0.58$ vs.\
LETKF's $+0.11$; for $\geq$15 components (40~episodes), the pattern
reverses (WS~$-0.44$ vs.\ LETKF~$+0.20$).  WaveSplat excels when the
spectral basis can resolve individual components, while LETKF's
grid-based formulation provides consistent but modest skill regardless
of spectral complexity.

\paragraph{LETKF variance deficit.}
LETKF's 100-episode VR of 0.127 is notably lower than the unlocalized
ETKF's 1.048, despite localization improving nRMSE and SS.
Gaspari--Cohn tapering at $r_{\text{loc}}=50$~m suppresses long-range
covariance, which reduces spurious spatial correlations but also
aggressively dampens ensemble spread in regions far from observations.
The resulting field has accurate spatial structure ($\rho=0.357$,
SS~$=+0.144$) but under-estimated amplitude---a mirror image of
ETKF's failure mode (correct amplitude, wrong structure).
This trade-off is inherent to localization radius selection and not
a tuning error: increasing $r_{\text{loc}}$ to 200~m raises VR to
$\sim$0.3 but degrades nRMSE to 0.279 (see localization sweep below).

\paragraph{Directional $R^2$ (supplementary diagnostic).}
Under the corrected tracker the 100-episode supplementary directional
$R^2$ is $+0.381$ (see Section~\ref{sec:dirdeficit}), the best among the
methods compared, but the aggregate is a mixture: $+0.983$ for
$\leq$5-component episodes against $-0.957$ for the 30-component ones,
where the beamforming initializer's peak-separation radius
(Sec.~\ref{sec:limitations}) leaves candidate peaks closely spaced and the
spectrum-space $R^2$ statistic becomes catastrophically sensitive to
single-bin peak misalignment.  The earlier draft reported this aggregate as
$-0.025$ with $-1.66$ on the 30-component subset; both figures predate the
tracker correction and are superseded.  The qualitative conclusion is
unchanged and now better supported: the spectral basis \emph{does} capture
directional structure when individual modes are resolvable, and loses it
when they are not.

\paragraph{Statistical significance.}
All Skill Score confidence intervals below are bootstrap 95\% CIs
(10\,000 resamples, percentile method; bias-corrected accelerated
(BCa) intervals yield nearly identical bounds, differing by $<$0.01
on all reported CIs).
On the full 100-episode dataset, both LETKF's SS~CI
($[{+0.13},\,{+0.16}]$) and WaveSplat's aggregate CI
($[{+0.06},\,{+0.28}]$) exclude zero: on the Nyquist-adequate
regenerated dataset WaveSplat's aggregate skill score is marginally but
significantly positive, whereas on the original coarse-mesh evaluation its
aggregate CI had included zero ($[{-0.08},\,{+0.15}]$).  The regime
dependence remains: restricting to $\leq$5-component
episodes (60/100), WaveSplat's SS~CI is $[{+0.50},\,{+0.65}]$, firmly
excluding zero and substantially exceeding LETKF's corresponding
$[{+0.09},\,{+0.13}]$.  For $\geq$15-component episodes the pattern
reverses (WS CI $[{-0.52},\,{-0.35}]$ vs.\ LETKF $[{+0.18},\,{+0.22}]$).
This complementarity suggests that an operational system could benefit
from complexity-aware method selection: WaveSplat for resolved seas
($\leq$5 components) and LETKF as a fallback for complex multi-modal
conditions.

A sensitivity sweep over localization radius ($r_{\text{loc}} \in
\{20, 50, 100, 200\}$~m) confirms that the reported LETKF performance
is near-optimal: $r_{\text{loc}}=20$~m and 50~m achieve nearly identical
nRMSE (0.234 vs.\ 0.234), while $r_{\text{loc}}=200$~m degrades
substantially (nRMSE 0.279, SS $-0.233$), validating that localization
is indeed beneficial and the chosen radius is reasonable.

\paragraph{Scalability to higher sea states.}
All experiments in this work are restricted to Beaufort~2--5
($H_s \leq 3$~m), where the linear superposition model introduces
$<0.5\%$ error.  At Beaufort~6--7 ($H_s = 3$--6~m), second-order
nonlinearities grow to $\sim$3--5\% and intermittent wave breaking
dissipates energy not captured by our forward model.  This limitation
is \emph{structural} for WaveSplat: the forward model
(Eq.~\ref{eq:forward}) is a linear superposition of plane waves, so
extending to higher sea states requires replacing the basis functions
(e.g., second-order Stokes components).  Grid-based methods such as
LETKF have a natural advantage here, as their forecast step can be
replaced by a nonlinear propagation model (e.g., HOS or NLSWE) without
modifying the data assimilation framework.

To verify this expectation, we generated 8 episodes at BF6
($H_s{=}3.0$~m) and BF7 ($H_s{=}5.0$~m) with varying complexity
(1--30 components) and ran both methods.  Under the linear simulator,
the same regime-dependent pattern persists:
WaveSplat dominates at ${\leq}5$ components
(mean SS~$+0.39$ vs LETKF $+0.16$) but degrades at ${\geq}15$
components (WS SS~$-0.28$ vs LETKF $+0.13$), with overall BF6--7
averages of nRMSE~0.212 (WS) vs 0.242 (LETKF) and SS~$+0.057$ (WS) vs
$+0.147$ (LETKF).
Since the simulator uses linear superposition throughout, these results
confirm that WaveSplat's performance depends on complexity rather than
sea state per se; \emph{real} BF6--7 performance would additionally
suffer from nonlinear model mismatch, likely favouring LETKF further.

In practice, Beaufort~2--5 covers the majority of autonomous vessel
operational hours---autonomous surface vessels typically initiate
storm-avoidance manoeuvres above Beaufort~5--6---but users targeting
heavy-weather applications should consider ensemble-based alternatives.

\subsection{Spectral parameter estimation accuracy}
\label{sec:spectral_accuracy}

A natural consequence of the grid-free representation is that it yields
individual spectral components---each with physically interpretable
wavenumber $(k_x, k_y)$, frequency $\omega$, amplitude $w$, and phase
$\phi$---which grid-based methods such as ETKF do not directly provide.
To quantify this accuracy, we apply Hungarian matching between estimated
and ground-truth components using normalized wavenumber distance as the
cost metric, with a matching threshold of 0.5.

\begin{table}[t]
\centering
\caption{Spectral parameter estimation accuracy by complexity, on the
Nyquist-adequate regenerated dataset (per-episode median over 5 estimator
seeds).
Detection rate: fraction of GT components matched.
Energy detection: fraction of total GT energy ($\sum A_i^2$) captured
by matched components.
$|k|$ error and direction error are averaged over matched pairs.}
\label{tab:spectral_accuracy}
\begin{tabular}{rcccc}
\toprule
$N_{\text{comp}}$ & Det.\ rate & Energy det. & $|k|$ rel.\ err. & Dir.\ err.\ [$^\circ$] \\
\midrule
1  & 100\% & 100\%  & 0.000 & 0.0 \\
2  & 100\% & 100\%  & 0.001 & 0.0 \\
5  & 100\% & 100\%  & 0.001 & 0.1 \\
15 &  60\% &  74.7\%         & 0.553 & 30.7 \\
30 &  33\% &  42.1\%         & 0.102 & 11.4 \\
\bottomrule
\end{tabular}
\end{table}

Table~\ref{tab:spectral_accuracy} reveals a clear separation by
spectral complexity.
For $\leq$5 components, every ground-truth component is recovered
(detection and energy-weighted detection both reach 100\%), with
wavenumber magnitude error below 0.1\% and directional error at or
below 0.1$^\circ$---from only $\sim$50 LiDAR points per frame.  On the
original coarse-mesh render this detection rate had sat at 50--60\%
because the shorter components fell below the reference mesh's spatial
Nyquist limit and were therefore unmatchable; the Nyquist-adequate
dataset recovers them, so the earlier ``high energy-detection despite
low raw-detection'' reading for $\leq$5-component seas was itself a
rendering artifact.
For $\geq$15 components, detection drops to 60--33\%, consistent with
the beamforming peak-separation radius discussed in
Sec.~\ref{sec:limitations}, which bounds but does not fully explain the
crowded-regime penalty.

Amplitude estimation, in contrast, is significantly degraded (mean
relative error 0.24--0.75, median over 5 estimator seeds; measured before
the tracker phase-advance correction and not re-measured, so read as an
upper bound---the aggregate variance ratio improved from $0.83$ to
$0.916$ under the correction).  The structural asymmetry behind it is
independent of that correction: wavenumber $(k_x, k_y)$ is determined
during initialization and frozen during tracking, while amplitude $w$ is
re-estimated every frame from sparse observations, so per-component
amplitude is the parameter with the least support.  We no longer
attribute the residual gap to a specific optimizer mechanism
(Section~\ref{sec:limitations}).  For downstream applications
such as wave-preview feedforward control, accurate wavenumber and
frequency estimation is arguably more important than amplitude, since
the controller primarily needs the correct \emph{timing and direction}
of incoming waves; amplitude can be independently estimated from
bulk statistics (e.g., $H_s$ from spectral moments).

\subsection{Sparse observation robustness: WaveSplat vs.\ LETKF}
\label{sec:sparse_comparison}

To quantify robustness to extreme observation sparsity, we sweep the
number of LiDAR points from 10 to 100 on four episodes and compare
WaveSplat against LETKF (Table~\ref{tab:sparse_comparison}).

\begin{table}[t]
\centering
\caption{Average nRMSE and Skill Score at varying observation counts
(4 episodes: BF3~2c, BF4~1c/5c/15c). WS~wins: episodes where
WaveSplat achieves lower nRMSE than LETKF.}
\label{tab:sparse_comparison}
\begin{tabular}{rccccr}
\toprule
$n_{\text{pts}}$ & WS nRMSE & LK nRMSE & WS SS & LK SS & WS wins \\
\midrule
10  & \textbf{0.191} & 0.239 & $+$0.293 & $+$0.090 & 3/4 \\
20  & \textbf{0.176} & 0.237 & $+$0.378 & $+$0.105 & 3/4 \\
30  & \textbf{0.180} & 0.235 & $+$0.340 & $+$0.117 & 3/4 \\
50  & \textbf{0.181} & 0.234 & $+$0.332 & $+$0.124 & 3/4 \\
100 & \textbf{0.181} & 0.234 & $+$0.327 & $+$0.124 & 3/4 \\
\bottomrule
\end{tabular}
\end{table}

WaveSplat maintains nRMSE~$\leq$0.191 even with only 10 observation
points, whereas LETKF nRMSE remains $\geq$0.234 regardless of
observation density.  This is because the spectral parameterization
($4N \leq 120$ unknowns) requires far fewer observations than the
grid discretization ($N_e \times G^2$ unknowns) to constrain the
solution.  In the sole exception (BF4~15c), LETKF outperforms
WaveSplat at all densities due to the high-complexity regime where
grid methods retain advantages (Section~\ref{sec:strongbaselines}).
For practical deployment, where LiDAR point counts often drop below
20 due to wave occlusion, spray, or fog, this sparse robustness is
critical.

